\medskip

Pour cet exercice, on admet que les deux suites $\left(a_n\right)$ et $\left(u_n\right)$ sont bien définies, ce qui équivaut à admettre que les termes de la suite $(u_n)$ sont tous différents de $-2$ (pour que la suite $\left(u_n\right)$ soit bien définie) et de $1$ (pour que $\left(a_n\right)$ soit bien définie).

\begin{enumerate}
\item On a : $u_{1} = \dfrac{2u_0 + 1}{ u_0 + 2} = \dfrac{2\times 2 + 1}{2 + 2 } = \dfrac{5}{4} = 1,25$.
\item 
	\begin{enumerate}
		\item On a : $a_{0} = \dfrac{u_0}{u_0 -1} = \dfrac{2}{2 - 1} = 2$

et $a_{1} = \dfrac{u_1}{u_1 - 1} = \dfrac{1,25}{1,25-1} = \dfrac{1,25}{0,25} = 5$.
		\item Soit $n$ un entier naturel.

On a, d'une part :\quad
	$\begin{aligned}[t]
a_{n+1} &= \dfrac{u_{n+1}}{u_{n+1}-1} = \dfrac{\dfrac{2u_n + 1}{u_n + 2}}{\dfrac{2u_n + 1}{u_n + 2} - 1} = \dfrac{\dfrac{2u_n + 1}{u_n + 2}}{\dfrac{2u_n + 1}{u_n + 2} - \dfrac{u_n + 2}{u_n + 2}}\\
	&=\dfrac{\dfrac{2u_n + 1}{u_n + 2}}{\dfrac{2u_n + 1- u_n - 2}{u_n + 2}} = \dfrac{2u_n + 1}{u_n + 2} \times \dfrac{u_n+2}{u_n -1}\\
	&=\dfrac{2u_n + 1}{u_n-1}
	\end{aligned}$

et, d'autre part :\quad
	$\begin{aligned}[t]
	3a_{n}-1 &= 3\times \dfrac{u_n}{u_n-1} - 1 = \dfrac{3 u_n}{u_n-1} - \dfrac{u_n-1}{u_n-1} = \dfrac{3u_n - (u_n - 1)}{u_n-1}\\
	&=\dfrac{2u_n + 1}{u_n-1}
	\end{aligned}$

On constate donc bien que, pour tout entier naturel $n$, $a_{n+1}=3 a_{n}-1$.

\textbf{Autre méthode :}

Soit $n$ un entier naturel.

$\begin{aligned}[t]
a_{n+1} & = \dfrac{u_{n+1}}{u_{n+1}-1}  =  \dfrac{u_{n+1}-1+1}{u_{n+1}-1}
= 1+ \dfrac{1}{u_{n+1}-1}\\
& = 1 + \dfrac{1}{\dfrac{2u_n + 1}{u_n + 2} - \dfrac{u_n + 2}{u_n + 2}} 
= 1 + \dfrac{1}{\dfrac{2u_n + 1}{u_n + 2}}\\
	& = 2+ \dfrac{u_n+2}{u_n -1} -1 \\
	& = \dfrac{2 ( u_n-1) + u_n+2}{u_n -1} -1 
	= \dfrac{2 u_n-2 + u_n+2}{u_n -1} -1\\
	& = \dfrac{3u_n}{u_n-1} - 1\\
	& = 3a_{n}-1
\end{aligned}$

\item Posons, pour tout entier naturel $n$ supérieur ou égal à 1, l'affirmation $P_n$ :

\hfill~\og $a_{n} \geqslant 3n - 1$\fg{}.\hfill

\emph{Initialisation :} On a calculé $a_1 = 5$ et donc on a bien $5 \geqslant 3\times 1 - 1 = 2$.

L'affirmation $P_1$ est donc vraie.

\emph{Hérédité :} Soit $n$ un entier naturel non nul. On suppose que l'affirmation $P_n$ est vraie, c'est-à-dire que $a_n \geqslant 3n-1$.

Par hypothèse de récurrence, on a :

$ \begin{aligned}[t]
a_n \geqslant 3n-1
	&\implies 3a_n \geqslant 3(3n-1) \quad \text{car } 3 > 0 \\
	&\implies 3a_n -1 \geqslant 9n - 3 - 1\\
	&\implies a_{n+1} \geqslant 9n-4 \quad \text{ d'après la question \textbf{2. b.} } 3an -1 = a_{n+1}\\
	&\implies a_{n+1} \geqslant 3 n + 6n + 3 - 7\\
	&\implies a_{n+1} \geqslant 3 (n+1) + 6k - 7\\
	&\implies a_{n+1} \geqslant 3 (n+1) - 1 + 6n - 6\\
	&\implies a_{n+1} \geqslant 3 (n+1) - 1 + 6(n-1)\\
	&\implies a_{n+1} \geqslant 3 (n+1) - 1 \quad \text{car } 6(n-1)\geqslant0\\
\end{aligned}$

Si, pour un entier $n$ naturel non nul (ce qui garantit $(n-1) \geqslant 0$), $P_n$ est vraie, alors $P_{n+1}$ est vraie également.

\emph{Conclusion :} L'affirmation $P_1$ est vraie, et, pour tout entier naturel non nul $n$, la véracité de l'affirmation est héréditaire : d'après le principe de récurrence, on en déduit que pour tout entier naturel non nul $n$, on a : $a_n \geqslant 3n - 1$.

\item Comme $3>0$, on a : $\lim\limits_{n \to +\infty} 3n - 1 = +\infty$.

Par comparaison, puisque pour tout $n$ non nul, on a $a_n \geqslant 3n-1$, on en déduit que la suite $(a_n)$ diverge vers $+\infty$.
\end{enumerate}

\item 
	\begin{enumerate}
		\item Soit $n$ un entier naturel.

$\begin{aligned}
a_n = \dfrac{u_n}{u_n - 1}
&\iff a_n \times (u_n - 1) = u_n \quad \text{car } u_n - 1 \neq 0 \\
&\iff a_n \times u_n - a_n - u_n = 0\\
&\iff (a_n - 1) \times u_n - a_n = 0\\
&\iff (a_n - 1) \times u_n = a_n\\
&\iff u_{n}=\dfrac{a_{n}}{a_{n}-1}
\end{aligned}$

Pour la dernière étape, la division par $a_n - 1$ est légitime, puisque $a_n$ est un quotient de deux nombres différents, car le dénominateur est égal au numérateur moins 1, donc $a_n$ ne peut pas être égal à 1, donc $a_n - 1$ est non nul.

Ainsi, on a bien exprimé $u_n$ en fonction de $a_n$ avec la relation annoncée dans le question.

		\item Avec les questions \textbf{2. a.} on a $a_0 = 2$ et avec \textbf{2. c.}, pour $n$ naturel non nul, on a $a_n \geqslant 3n - 1 > 0$.

Ainsi, pour tout entier naturel $n$, $a_n$ est non nul.

On déduit donc de la question précédente que, pour tout $n$ naturel, on a : $u_n = \dfrac{a_n}{a_n -1} = \dfrac{1}{1 - \dfrac{1}{a_n}}$.

Enfin, puisque $(a_n)$ diverge vers $+\infty$, par limite du quotient, on a : $\lim\limits_{n \to +\infty}~\dfrac{1}{a_n}~=~0$,

puis, par limite de la somme $\lim\limits_{n \to +\infty} 1 - \dfrac{1}{a_n} = 1$,

enfin, par limite du quotient : $\lim\limits_{n\to +\infty} u_n = \lim\limits_{n\to +\infty}\dfrac{1}{1-\dfrac{1}{a_n}} = 1$.

La suite $\left(u_{n}\right)$ converge donc vers 1.
	\end{enumerate}
\item Puisqu'on admet que la suite $\left(u_{n}\right)$ est décroissante, cela signifie que la suite est minorée par sa limite et donc que pour tout entier naturel $n$, on a $u_n \geqslant 1$ donc $u_n - 1 \geqslant 0$.
	\begin{enumerate}
		\item Cette fonction python initialise la variable \texttt{u} à 2, ce qui est la valeur de $u_0$ et la variable \texttt{n} à 0, ce qui est l'indice correspondant.

Puis, à chaque exécution de la boucle \texttt{while},  \texttt{u} se voit affecter le terme suivant dans la suite $(u_n)$ et \texttt{n} se voit affecter l'entier suivant, donc l'indice correspondant au terme dont la valeur est stockée dans la variable \texttt{u}.

Cette boucle \texttt{while} tourne tant que l'écart entre le terme stocké dans la variable \texttt{u} et la limite 1 est strictement supérieur à la valeur \texttt{p}, qui est l'argument choisi pour invoquer la fonction.

Les valeurs \texttt{n} et \texttt{u} renvoyées par l'appel de la fonction \texttt{algo(p)} correspondent donc respectivement à l'indice et à la valeur du premier terme de la suite pour lequel l'écart entre le terme et la limite de la suite est inférieur ou égal à la valeur \texttt{p} choisie.

		\item On parcourt la suite à la calculatrice, et on constate que $u_5 \approx \np{1,0027}$, donc $u_5 - 1 > 0,001$ et $u_6 \approx \np{1,0009}$, donc $u_6 - 1 \leqslant 0,001$.

La valeur de $n$ pour $p=0,001$ est donc 6 (et la valeur renvoyée pour \texttt{u} est donc une valeur approchée de $u_6$).

On peut aussi programmer la fonction python sur la calculatrice et faire l'appel \texttt{algo(0.001)}, qui renvoie \texttt{(6, \np{1.000914913083257})}
	\end{enumerate}
\end{enumerate}

\bigskip
